% ============================================================================
\chapter{Wiles 的证明策略}
\label{ch:wiles-proof}
% ============================================================================

% ----------------------------------------------------------------------------
\section{证明的整体架构}
% ----------------------------------------------------------------------------

% TODO: 
% Wiles 证明模性猜想（半稳定情形）的路线图：
% 1. 将问题归结为 ρ̄_{E,3} 的模性提升
% 2. R = T 定理
% 3. Taylor-Wiles 方法

% ----------------------------------------------------------------------------
\section{残余表示与模性提升}
% ----------------------------------------------------------------------------

% TODO: 
% - 残余表示 ρ̄: G_Q → GL₂(F_ℓ)
% - Serre 猜想的背景
% - 从 ρ̄ 的模性到 ρ 的模性（提升问题）

% ----------------------------------------------------------------------------
\section{变形环与 Hecke 代数}
% ----------------------------------------------------------------------------

% TODO:
% - 万有变形环 R
% - Hecke 代数 T
% - R = T 定理的陈述

% ----------------------------------------------------------------------------
\section{Taylor-Wiles 方法}
% ----------------------------------------------------------------------------

% TODO:
% - 辅助素数集（Taylor-Wiles primes）
% - 极限论证
% - R = T 的证明思路

% ----------------------------------------------------------------------------
\section{Langlands-Tunnell 定理与 3-5 互换}
% ----------------------------------------------------------------------------

% TODO:
% - 用 Langlands-Tunnell 启动 ℓ = 3 的情形
% - 当 ρ̄_{E,3} 不适用时切换到 ℓ = 5

% ----------------------------------------------------------------------------
\section{习题}
% ----------------------------------------------------------------------------
